\section{The survival amplitude \texorpdfstring{\(A(p)\)}{A(p)}}
\label{sec:constant-A}

For \(0<p<1/2\), all of the limiting conditional moments obtained in
\cref{sec:qsd} depend on the amplitude
\begin{equation}
  A(p)=\lim_{n\to\infty}\frac{S_n}{(2p)^n}.
  \label{eq:A-def-repeat}
\end{equation}
Direct iteration gives accurate values away from criticality, but it does not
explain the constant and becomes slow as \(p\uparrow1/2\).  We now derive an exact
product, an exact series, two-sided bounds, a rigorous near-critical asymptotic and
a structural connection with linearisation in complex dynamics.  These results are
specific to the binary process and are therefore proved in full.

\subsection{An infinite product}
\label{sec:A-product}

Put
\begin{equation}
  r=2p\in(0,1),
  \qquad
  w_n=\frac{S_n}{2}.
\end{equation}
Then the survival recursion becomes the logistic map
\begin{equation}
  w_{n+1}=f_r(w_n):=rw_n(1-w_n),
  \qquad w_0=\frac12.
  \label{eq:logistic-map}
\end{equation}
Only the attracting regime \(0<r<1\) is relevant here.  Define
\(x_n=w_n/r^n\).  Successive ratios telescope:
\begin{equation}
  \frac{x_{n+1}}{x_n}
  =\frac{w_{n+1}}{rw_n}=1-w_n.
  \label{eq:x-ratio}
\end{equation}
Since \(A(p)=2\lim_nx_n\) and \(2x_0=1\), this suggests the following exact
representation.

\begin{proposition}[Product representation]
\label{prop:A-product}
For \(0<p<1/2\),
\begin{equation}
  A(p)=\prod_{k=0}^{\infty}(1-w_k)
  =\frac12\prod_{k=1}^{\infty}\left(1-\frac{S_k}{2}\right).
  \label{eq:A-product}
\end{equation}
The product converges to a strictly positive number.  In particular, the limit in
\cref{eq:A-def-repeat} exists.
\end{proposition}

\begin{proof}
Iteration of \cref{eq:x-ratio} gives
\begin{equation}
  x_n=x_0\prod_{k=0}^{n-1}(1-w_k).
\end{equation}
Moreover, \(0<w_{n+1}<rw_n\), so \(w_n\leq w_0r^n\).  Hence
\(\sum_nw_n<\infty\).  An infinite product with factors \(1-w_n\in(0,1)\) has a
positive limit whenever \(\sum_nw_n\) converges.  Passing to the limit and using
\(w_0=1/2\) proves both identities in \cref{eq:A-product}.
\end{proof}

The finite product
\begin{equation}
  A_N(p)=\frac12\prod_{k=1}^{N}\left(1-\frac{S_k}{2}\right)
  \label{eq:A-partial-product}
\end{equation}
decreases to \(A(p)\), so every truncation is an upper bound.  A useful lower bound
comes from \(\prod_k(1-a_k)\geq1-\sum_ka_k\) for \(a_k\in[0,1]\).  Because
\(w_{k+1}\leq rw_k\),
\begin{equation}
  A_N\left(1-\frac{w_{N+1}}{1-r}\right)
  \leq A(p)\leq A_N,
  \label{eq:A-product-tail-bound}
\end{equation}
whenever the left factor is positive.  The bracket provides a stopping criterion
that is stronger than merely observing that successive floating-point iterates have
stopped changing.

\subsection{An exact series and two-sided bounds}
\label{sec:A-series-bounds}

The product proves existence, but an additive representation is more effective near
criticality.  Set
\begin{equation}
  v_n=\frac{r^n}{S_n}.
\end{equation}
Then \(v_n\to1/A(p)\), and the factorised recursion
\(S_{n+1}=rS_n(1-S_n/2)\) gives
\begin{align}
  v_{n+1}-v_n
  &=\frac{r^{n+1}}{rS_n(1-S_n/2)}-\frac{r^n}{S_n}
    \notag\\
  &=\frac{r^n}{2-S_n}.
  \label{eq:v-increment}
\end{align}
Since \(v_0=1\), summation yields an exact series.

\begin{proposition}[Series representation]
\label{prop:A-series}
For \(0<p<1/2\),
\begin{equation}
  \frac{1}{A(p)}
  =1+\sum_{n=0}^{\infty}\frac{(2p)^n}{2-S_n}.
  \label{eq:A-series}
\end{equation}
\end{proposition}

The nonlinearity is confined to the denominator, while the numerator is geometric.
Since \(0<S_n\leq1\), we have
\begin{equation}
  \frac12r^n<\frac{r^n}{2-S_n}\leq r^n,
  \label{eq:A-series-term-bounds}
\end{equation}
with strictness on the appropriate side for at least one term.  Write
\(\varepsilon=1-r=1-2p\).  Summing the geometric bounds and inverting gives the
following result.

\begin{proposition}[Two-sided bounds]
\label{prop:A-bounds}
For \(0<p<1/2\),
\begin{equation}
  \frac{\varepsilon}{1+\varepsilon}
  <A(p)<
  \frac{2\varepsilon}{1+2\varepsilon}.
  \label{eq:A-bounds-epsilon}
\end{equation}
Equivalently,
\begin{equation}
  \frac{1-2p}{2(1-p)}
  <A(p)<
  \frac{2(1-2p)}{3-4p}.
  \label{eq:A-bounds-p}
\end{equation}
The lower bound is one half of the continuous-time constant:
\begin{equation}
  A(p)>\frac12A_{\mathrm c}(p).
  \label{eq:A-half-Ac}
\end{equation}
\end{proposition}

\begin{proof}
By \cref{eq:A-series-term-bounds},
\begin{equation}
  1+\frac{1}{2\varepsilon}
  <\frac{1}{A(p)}
  <1+\frac{1}{\varepsilon}.
\end{equation}
All quantities are positive, so inversion reverses the inequalities and gives
\cref{eq:A-bounds-epsilon}.  Substitution of
\(\varepsilon=1-2p\) gives \cref{eq:A-bounds-p}.  Finally,
\(A_{\mathrm c}(p)=2\varepsilon/(1+\varepsilon)\) by
\cref{eq:Ac-definition}, proving \cref{eq:A-half-Ac}.
\end{proof}

There is also a bound that comes from the support of the conditional population.

\begin{proposition}[Parity bound]
\label{prop:A-parity}
For \(0<p<1/2\),
\begin{equation}
  A(p)\leq\frac12,
  \label{eq:A-parity}
\end{equation}
and \(A(p)\to1/2\) as \(p\downarrow0\).
\end{proposition}

\begin{proof}
For every \(n\geq1\), the binary process has even population size.  Conditional on
survival, \(Z_n\geq2\), so \cref{eq:conditional-mean-exact} gives
\(m^n/S_n\geq2\).  Letting \(n\to\infty\) yields \(1/A(p)\geq2\).
On the other hand, the lower bound in \cref{eq:A-bounds-p} tends to \(1/2\) as
\(p\downarrow0\); together with \cref{eq:A-parity}, this proves the endpoint limit.
\end{proof}

The definition is therefore extended continuously by
\begin{equation}
  A(0):=\lim_{p\downarrow0}A(p)=\frac12.
  \label{eq:A-at-zero}
\end{equation}
This convention avoids the undefined ratio \(S_n/(2p)^n\) at \(p=0\) itself.

\begin{table}[htbp]
\centering
\caption{High-precision values of \(A(p)\), the bounds in
\cref{prop:A-bounds}, the near-critical approximation \(2(1-2p)\), and the
conditional mean \(1/A(p)\).  The last column is the number of product factors used
by the reference calculation.  The values are reproduced by
\texttt{scripts/reproduce\_numerics.py}.}
\label{tab:A-values}
\small
\setlength{\tabcolsep}{4.2pt}
\begin{tabular}{@{}rrrrrrr@{}}
\toprule
\(p\) & \(A_{\mathrm c}/2\) & \(A(p)\) & upper bound & \(2(1-2p)\) & \(1/A(p)\) & factors\\
\midrule
0.05   & 0.473684 & 0.486180 & 0.642857 & 1.800000 & 2.0568 & 14\\
0.10   & 0.444444 & 0.469382 & 0.615385 & 1.600000 & 2.1305 & 20\\
0.20   & 0.375000 & 0.423895 & 0.545455 & 1.200000 & 2.3591 & 35\\
0.30   & 0.285714 & 0.353967 & 0.444444 & 0.800000 & 2.8251 & 62\\
0.40   & 0.166667 & 0.237647 & 0.285714 & 0.400000 & 4.2079 & 143\\
0.45   & 0.090909 & 0.145069 & 0.166667 & 0.200000 & 6.8933 & 303\\
0.48   & 0.038462 & 0.068011 & 0.074074 & 0.080000 & 14.7036 & 786\\
0.49   & 0.019608 & 0.036395 & 0.038462 & 0.040000 & 27.4764 & 1\,591\\
0.499  & 0.001996 & 0.003944 & 0.003984 & 0.004000 & 253.519 & 16\,095\\
0.4999 & 0.000200 & 0.000399 & 0.000400 & 0.000400 & 2504.65 & 161\,156\\
\bottomrule
\end{tabular}
\end{table}

\subsection{The near-critical limit}
\label{sec:A-near-critical}

The bounds in \cref{eq:A-bounds-epsilon} show that \(A(p)\) is of order
\(\varepsilon\), but leave a factor of two unresolved.  The exact series determines
the leading constant and the order of the first relative correction.

\begin{theorem}[Near-critical survival amplitude]
\label{thm:A-near-critical}
As \(p\uparrow1/2\), with \(\varepsilon=1-2p\downarrow0\),
\begin{equation}
  A(p)=2\varepsilon
  \left[1+O\!\left(\varepsilon\log\frac1\varepsilon\right)\right].
  \label{eq:A-near-critical}
\end{equation}
In particular,
\begin{equation}
  A(p)\sim2(1-2p),
  \qquad
  \frac{1}{A(p)}\sim\frac{1}{2(1-2p)}.
  \label{eq:A-leading-asymptotic}
\end{equation}
\end{theorem}

\begin{proof}
Write \(r=1-\varepsilon\).  The identity
\begin{equation}
  \frac{1}{2-S_n}=\frac12+\frac{S_n}{2(2-S_n)}
\end{equation}
turns \cref{eq:A-series} into
\begin{equation}
  \frac1{A(p)}
  =1+\frac1{2\varepsilon}+D_\varepsilon,
  \qquad
  D_\varepsilon
  :=\sum_{n=0}^{\infty}
     r^n\frac{S_n}{2(2-S_n)}.
  \label{eq:A-series-decomposition}
\end{equation}
We require only a logarithmic bound on the remainder.  The map
\(s\mapsto ps(2-s)\) is increasing in both \(p\) and \(s\) on the relevant range,
so the subcritical survival probability is bounded above by its critical counterpart,
\(S_n(p)\leq S_n(1/2)\).  From \cref{eq:reciprocal-increment},
\begin{equation}
  \frac1{S_{n+1}(1/2)}-\frac1{S_n(1/2)}\geq\frac12,
\end{equation}
and hence \(S_n(1/2)\leq2/(n+2)\).  Since \(2-S_n\geq1\),
\begin{align}
  0\leq D_\varepsilon
  &\leq\frac12\sum_{n=0}^{\infty}r^nS_n(p)
   \leq\sum_{n=0}^{\infty}\frac{r^n}{n+2}
  \notag\\
  &=\frac{-\log(1-r)-r}{r^2}
   =O\!\left(\log\frac1\varepsilon\right).
  \label{eq:D-epsilon-bound}
\end{align}
Factoring \(1/(2\varepsilon)\) from \cref{eq:A-series-decomposition} gives
\begin{equation}
  \frac1{A(p)}
  =\frac1{2\varepsilon}
   \left[1+2\varepsilon\{1+D_\varepsilon\}\right].
\end{equation}
The quantity after the leading one is
\(O(\varepsilon\log(1/\varepsilon))\); inversion proves
\cref{eq:A-near-critical}.
\end{proof}

The proof locates the logarithm that slows convergence to the leading asymptotic.
Numerically, the ratios \(A/(2\varepsilon)\) are approximately
\(0.90987,0.98612,0.99814,0.99977,0.999972\) as \(\varepsilon\) decreases from
\(2\times10^{-2}\) to \(2\times10^{-6}\).  An exploratory fit gives
\begin{equation}
  1-\frac{A}{2\varepsilon}
  \approx\varepsilon\left(\log\frac1\varepsilon+0.78\right),
  \label{eq:A-empirical-correction}
\end{equation}
but the constant \(0.78\) is empirical and is not used subsequently.  The proved
statement is the order estimate in \cref{eq:A-near-critical}.

\subsection{Discrete and continuous time}
\label{sec:A-discrete-continuous}

The continuous-time amplitude from \cref{eq:Ac-definition} is
\begin{equation}
  A_{\mathrm c}(p)=\frac{1-2p}{1-p}
  =\frac{2\varepsilon}{1+\varepsilon}.
  \label{eq:Ac-epsilon}
\end{equation}
The lower bound in \cref{eq:A-bounds-epsilon} is exactly
\(A_{\mathrm c}/2\).  At the strongly subcritical endpoint,
\begin{equation}
  \frac{A(p)}{A_{\mathrm c}(p)}\longrightarrow\frac12
  \qquad(p\downarrow0),
  \label{eq:A-ratio-zero}
\end{equation}
whereas \cref{eq:A-leading-asymptotic,eq:Ac-epsilon} give
\begin{equation}
  \frac{A(p)}{A_{\mathrm c}(p)}\longrightarrow1
  \qquad(p\uparrow\tfrac12).
  \label{eq:A-ratio-critical}
\end{equation}
The factor of two at \(p\downarrow0\) is a support effect.  A surviving binary
generation contains at least two individuals, while a surviving continuous-time
population may contain one.  The respective limiting conditional means are therefore
two and one.  Near criticality this distinction is lower order, and both amplitudes
share the leading behaviour \(2\varepsilon\).

Computed values of \(A/A_{\mathrm c}\) increase from approximately
\(0.503,0.528,0.619,0.798,0.928\) to \(0.988\) at
\(p=0.01,0.1,0.3,0.45,0.49,0.499\), respectively.  This numerical monotonicity is
descriptive; no monotonicity theorem is required for
\cref{eq:A-ratio-zero,eq:A-ratio-critical}.

\begin{figure}[htbp]
  \centering
  \includegraphics[width=0.54\textwidth]{dtctA.png}
  \caption{The elementary continuous-time amplitude \(A_{\mathrm c}(p)\) and
  numerical values of the discrete-time \(A(p)\).  The additional black curve is
  an exploratory symbolic-regression approximation discussed in
  \cref{app:closed-form-search}; it is visually close but fails the critical
  asymptotic.}
  \label{fig:A-discrete-continuous}
\end{figure}

\subsection{Koenigs linearisation}
\label{sec:A-koenigs}

The product in \cref{eq:A-product} is also the classical linearising coordinate of
the logistic map.  For \(f_r(z)=rz(1-z)\), \(0<r<1\), Koenigs' theorem
\cite{Koenigs1884} gives a unique function \(\psi_r\), analytic near the attracting
fixed point \(0\), such that
\begin{equation}
  \psi_r(f_r(z))=r\psi_r(z),
  \qquad
  \psi_r(0)=0,
  \qquad
  \psi_r'(0)=1.
  \label{eq:schroeder-koenigs}
\end{equation}
The function extends along the attracting basin by the functional equation.  In the
\(\psi_r\)-coordinate, one nonlinear iteration of \(f_r\) is simply multiplication
by \(r\).

\begin{figure}[htbp]
  \centering
  \includegraphics[width=0.70\textwidth]{figure2_koenigs_linearization.png}
  \caption{Koenigs linearisation.  The coordinate change \(\psi_r\) conjugates the
  nonlinear map \(f_r\) to multiplication by \(r\), expressing
  \cref{eq:schroeder-koenigs}.}
  \label{fig:koenigs-linearisation}
\end{figure}

\begin{proposition}[The survival amplitude as a Koenigs value]
\label{prop:A-koenigs}
For \(0<p<1/2\) and \(r=2p\),
\begin{equation}
  A(p)=2\psi_r\!\left(\frac12\right).
  \label{eq:A-koenigs}
\end{equation}
\end{proposition}

\begin{proof}
Let \(w_{n+1}=f_r(w_n)\) with \(w_0=1/2\).  Iterating
\cref{eq:schroeder-koenigs} gives
\begin{equation}
  \psi_r(w_0)=\frac{\psi_r(w_n)}{r^n}.
  \label{eq:koenigs-orbit}
\end{equation}
Since \(w_n\to0\) and \(\psi_r(z)/z\to1\) as \(z\to0\),
\begin{equation}
  \psi_r(w_0)=\lim_{n\to\infty}\frac{w_n}{r^n}.
\end{equation}
But \(w_n=S_n/2\), so multiplication by two gives \cref{eq:A-koenigs}.
\end{proof}

This identification explains why elementary curve fitting is an unreliable route to
\(A(p)\).  Becker and Bergweiler's classification
\cite{BeckerBergweiler1995} shows that Schröder conjugacies of rational maps are
hypertranscendental outside a short list of exceptional dynamical families.  The
quadratic maps corresponding to \(r\in(0,1)\) are not exceptional; the conjugacy
calculation is recorded in \cref{app:quadratic-conjugacy}.  Thus, for each fixed
\(r\in(0,1)\), the function \(z\mapsto\psi_r(z)\) satisfies no algebraic
differential equation over the rational functions and, in particular, is not
elementary.

The conclusion must stop there.  Hypertranscendence in the argument \(z\) does not
prove that the parameter map
\begin{equation}
  p\longmapsto2\psi_{2p}(1/2)=A(p)
  \label{eq:A-parameter-map}
\end{equation}
is non-elementary, because this map varies \(r\) while fixing the argument.  No
elementary expression for \(A(p)\) is known, and the Koenigs representation gives a
structural reason not to expect one, but it is not a transcendence theorem for the
values in \cref{eq:A-parameter-map}.  The exploratory searches and the conjugacy to
the quadratic family are recorded, with this limitation, in
\cref{app:closed-form-search}.

\subsection{Computation in practice}
\label{sec:A-computation}

Away from criticality, the product \cref{eq:A-partial-product} is direct, monotone
and accompanied by the certified bracket \cref{eq:A-product-tail-bound}.  Near
criticality the number of factors grows on the scale \(1/(1-2p)\), with an additional
logarithmic cost if many digits are required.  The final column of
\cref{tab:A-values} shows this slowdown.

For calculations that do not require a certified product value arbitrarily close to
criticality, the exact product and the asymptotic can be combined:
\begin{equation}
  A(p)\approx
  \begin{cases}
    \displaystyle
    \frac12\prod_{k=1}^{N}\left(1-\frac{S_k}{2}\right),
      &p<c,\\[4mm]
    2(1-2p),&c\leq p<\tfrac12.
  \end{cases}
  \label{eq:A-practical}
\end{equation}
Here \(N\) is chosen from \cref{eq:A-product-tail-bound}, and \(c\) is selected to
match the required accuracy.  The relative error of the second branch is
\(O(\{1-2p\}\log\{1/(1-2p)\})\) by \cref{thm:A-near-critical}.  For example,
Choosing \(c=0.4999\) gives a sub-percent approximation across the whole interval;
\(c=0.499\) leaves an error of about \(1.4\%\) at the join.  A smooth blend can be
used if a differentiable surrogate is needed.

All values in this chapter are generated from the product rather than from a fitted
formula.  The accompanying script uses decimal arithmetic, reports the rigorous tail
bracket and switches to the asymptotic only when explicitly requested.

\FloatBarrier
